\SourceLeaf{155}
deren Studium erst für denjenigen von lebhaftem Interesse sein kann, der
die Methoden bereits anzuwenden versteht. Ich pflege drei Begründungsarten
vorzutragen.

\begin{enumerate}
\item[1)] eine, die allein auf Zweckmäßigkeitsüberlegungen beruht und sich leicht und sehr
einleuchtend darlegen lässt,
\item[2)] die in der Th. M. C. C. gelehrte Art der Verbindung mit der
Wahrscheinlichkeitsrechnung,
\item[3)] die davon \textls{völlig verschiedene}, in der Theoria Combinationis
Observationum dargelegte Art der Verbindung mit der Wahrscheinlichkeitsrechnung.
Sie ist nach meiner Überzeugung die einzig
zulässige.
\end{enumerate}

Für jeden, dem diese Lehre noch ganz unbekannt ist, halte ich diese
Reihenfolge für die zweckmäßigste. Darstellungen anderer kenne ich aber
keine, die ich empfehlen könnte. Was aber No. 1) betrifft, so kann Herr
\textsc{Warensdorf} dies aus unzähligen Büchern lernen (soweit es um das A. B. C. der Sache
geht), und es ist ziemlich gleichgültig, welches er wählt, sofern er später die Gelegenheit nutzt,
für sich zu verbessern, was fast alle schlecht darstellen.
So betonen diese Bücher zum Beispiel nicht hinreichend,
dass man auch dann, wenn die beobachteten Größen bereits lineare
Funktionen der unbekannten Elemente sind (wie bei der Pendellänge
$x+y\sin\varphi^2$), nicht diese Elemente als Unbekannte der Aufgabe betrachten soll.
Vielmehr soll man die Korrekturen wählen, die an bereits bekannten
\textls{möglichst guten Näherungswerten} anzubringen sind, und so weiter.

. . . . . . . . . . . .
